\documentclass[10pt]{beamer}
\usepackage{csquotes}% Recommended

\usepackage{graphicx}
\usepackage{xcolor}
\usepackage{amsfonts}
\usepackage{amsmath}
\usepackage{amssymb}
\usepackage[skip=10pt plus1pt] {parskip}
\usepackage{multicol}
\usepackage{multirow}

%Information to be included in the title page:
\title{Continuous Probability}
\author{UChicago Political Science Math Camp}
\institute{Ryan Gibbons \& Luke Cain}
\date{2025}

\begin{document}

\frame{\titlepage}

\begin{frame}{Continuous Probability}

We can expand our understanding of probability to include infinite outcomes \\

\vspace{.5cm}

We know that a random variable can model discrete outcomes: i.e., flipping a coin, rolling a dice \\

\vspace{.5cm}

What if we took a dice with one face for every real number $\mathbb{R} \in [0,1]$?
    
\end{frame}

\begin{frame}{PDFs}
    A PDF extends a PMF to include a continuous domain, such that

    \[f(x) = \Pr[X=x]\]

    Where $f$ returns a real number for every input of $x$. A PDF returns the probability density of a certain outcome of a random variable.
\end{frame}

\begin{frame}{Probability Axioms}
    Recall our probability axioms:

    \begin{enumerate}
        \item $\Pr[a] \geq 0 \; \forall \; a \in A$
        \item $\sum_{a \in A} \Pr[a] = 1$
        \item For disjoint outcomes, $\Pr[\bigcup_{i=1}^k a_i] = \sum_{i=1}^k \Pr[a_i]$
    \end{enumerate}

    How might we apply these to continuous probability distributions?
\end{frame}

\begin{frame}{PDF Warnings}
    A few things become questionable with PDFs:
    \begin{enumerate}
        \item PDFs do not need to be continuous nor differentiable.
        \item PDFs can be defined for negative inputs (RV outcomes), just not negative outputs (probabilities)
        \item PDFs do not need to be injective nor surjective
    \end{enumerate}

    \textit{Conceptualize: what is this PDF telling me?}
\end{frame}

\begin{frame}{PDF Bounds}
    It's not uncommon to see a PDF take a form such as:

    \[f(x) = 2x\]

    But, the condition $\int_{\mathbb{R}} f(x)dx = 1$ still must hold. We then create \textbf{bounds} which allow this: \pause

    \[\int_a^b 2xdx = 1 \implies x^2 \bigg|_a^b = 1\]
    \[\implies b^2 - a^2 = 1 \implies b = \sqrt{a^2+1} \]

    More generally, we must choose $a,b$ such that $F(b)-F(a) = 1$.
\end{frame}

\begin{frame}{CDFs}
    Say for some PDF $f(X)$, we want to know the likelihood that $x$ is larger than some $X$. We can use the \textbf{cumulative distribution function}:

    \[F(x) = \Pr[X \leq x] = \int_{-\infty}^X f(s) ds\]

    We use FTOC to obtain this: $\int_a^b f(x)dx = F(b) - F(a)$, for $a = - \infty$, and some arbitrary $s$
\end{frame}

\begin{frame}{Expanded CDFs}
    This FTOC definition allows us to specify the probability of an RV falling within a given range. For example,

    \[\Pr[a < x < b] = \int_a^b f(x)dx\] \pause

    We can then use our \textbf{disjoint probability axiom} to combine multiple domains:

    \[\Pr[x \in [a,b] \text{ OR } [c,d]] = \int_a^b f(x)dx + \int_c^d f(x)dx\]
\end{frame}

\begin{frame}{CDF Foundations}
    What do our first two axioms imply about continuous CDFs? \pause \\
\vspace{1cm}
    \begin{itemize}
        \item Weakly increasing: $F(x) \ge F(x') \; \forall \; x > x'$
        \item Comprehensive: $F(\infty) = 1$
    \end{itemize}
\end{frame}

\begin{frame}{Joint Distributions}
    Joint distributions are an easy extension from discrete probability:

    \[p(x,y) = P(X = x, Y = y)\]

    \[\implies f(x,y); P((X,Y) \in A) = \int \int_A f(x,y)dxdy\]

    \textit{Probability represented by volume under the joint PDF}
\end{frame}

\begin{frame}{Marginal Distributions}
    A \textbf{marginal distribution}, denoted $f_X(x)$, reflects the distribution of one variable, holding the other variable(s) constant. \\

    \vspace{1cm}

    Then, we can represent by integrating \textcolor{red}{w.r.t. the other variable!}

    \[f_X(x) = \int_{-\infty}^{\infty} f(x,y) dy\]
    \[f_Y(x) = \int_{-\infty}^\infty f(x,y) dx\]

    \[\]
\end{frame}

\begin{frame}{Joint PDFs}
    You cannot always derive a joint PDF, but: if two events are independent, their joint PDF is the product of their marginal PDFs \\

    \vspace{1cm}

    Example: Let $f(x) = e^{-x}$ and $f(y) = N(\bar{x},\sigma)$. What is $f(x,y)$?
\end{frame}

\begin{frame}{Conditional Probability}
    To evaluate the \textbf{conditional probability} of a continuous RV, we hold the other variable constant:

    \[\Pr[y | x] = f_{y|x} (y|x) = \frac{f(x,y)}{f_X(x)}\]
\end{frame}

\begin{frame}{Uniform Distribution}
    The uniform distribution takes a slightly different form for continuous probability:

    \[f(x) = \frac{1}{b-a}\]

    for uniformity over the domain $[a,b]$.
\end{frame}

\begin{frame}{Expanding Expectation}

Recall that for a discrete RV, $E[X] = \sum_{x \in X} \Pr[x] \cdot x$ \\

\vspace{.5cm}

How can we expand this to an infinite cardinality of $x$? \pause \textit{Would we need, perhaps, an infinite sum?} \pause

\[E[X] = \int_{-\infty}^{\infty} x f(x) dx\]

\textcolor{gray}{\textit{Unfortunately, you can't isolate out $\int_{\mathbb{R}} f(x)dx = 1$ here}}
\end{frame}

\begin{frame}{An Integration Corollary:}
    The sum rule lends that

    \[\int_a^c f(x)dx = \int_a^b f(x) dx + \int_b^c f(x) dx\]

    Suppose $f(x)$ is only defined on $[a,b]$ and therefore takes a value of $0$ for all other inputs. Then,

    \[\int_a^c f(x)dx = \int_a^b f(x)dx + \int_b^c 0dx\]

    \[\implies \int_a^b f(x)dx\]

    \pause \textcolor{red}{Much of probability calls for infinite integrals, but we can often isolate a well-defined domain.}
\end{frame}

\begin{frame}{Median}
    The median defines a value such that half of the mass is above that threshold, and half is below. \\

    \vspace{.5cm}

    How could we apply that to continuous calculus? \pause

    \[\text{Median}(X) \equiv x \; \text{ s.t. } \; \int^x f(X)dX = \frac{1}{2}\]
\end{frame}





\end{document}